Unless otherwise specified, the following terms are used throughout this document.

\begin{description}
    \item[\must] Indicates a mandatory requirement. A deviation requires an approved exception.

    \item[\mustnot] Indicates a mandatory prohibition. The described action or condition is not permitted unless an approved exception explicitly allows it.

    \item[\should] Indicates a strong recommendation. A deviation is permitted when there is a documented technical reason.

    \item[\shouldnot] Indicates a strong recommendation to avoid an action or condition. A deviation should have a documented technical reason.

    \item[\may] Indicates permission or an optional practice. It does not impose a requirement.

    \item[\donot] Indicates a direct instruction to avoid an action. It is normally used for procedural instructions and prohibitions. 

    \item[Project] Refers to the software, documentation, tools, configuration, infrastructure, and processes controlled by this document.

    \item[Source Code] Human readable code written in a programming language, including prod code, test code, scripts, and supporting utilities.

    \item[Build] The process of transforming source code, configuration, and other inputs into executable programs, libraries, packages, or other generated artifacts.

    \item[Development Environment] The hardware, operating system, software, tools, libraries, settings, and services used to develop or maintain the project.

    \item[Change] Any addition, deletion, or modification to project code, documentation, configuration or tooling.

    \item[Release] A formally identified version of the project made available for use or further integration.

    \item[Exception] An approved deviation from a mandatory requirement in this document.

    % \item[Configuration Item] A project artifact placed under version control or otherwise managed as part of the project's controlled baseline.
\end{description}

The words ``\must'', ``\should'', and ``\may'' are interpreted according to the definitions above. These terms are intentionally used consistently so that the strength of each rule is clear.